\begin{table}[h!]
    \scriptsize
    \captionsetup{width=\linewidth}
    \Caption{\label{tab:criterios_parada} Resumo dos criterios de parada na melhor execucao. Legenda: AT representa ACO\_R+TS e HR representa HHO+RVNS; Metodo identifica a metaheuristica; Crit. identifica o criterio de parada; $q$ e o numero de melhores execucoes encerradas por esse criterio. tempo indica limite de tempo; estag. indica estagnacao; iter. indica numero maximo de iteracoes.}
    \IBGEtab{}{
        \begin{tabular}{llll}
            \toprule
            Conj. & Metodo & Crit. & $q$ \\
            \midrule \midrule
            TODOS & HR & iter. & 6 \\
            TODOS & HR & estag. & 63 \\
            TODOS & HR & tempo & 61 \\
            TODOS & AT & iter. & 15 \\
            TODOS & AT & estag. & 112 \\
            TODOS & AT & tempo & 3 \\
            CUBIC & HR & iter. & 2 \\
            CUBIC & HR & estag. & 11 \\
            CUBIC & HR & tempo & 17 \\
            CUBIC & AT & iter. & 8 \\
            CUBIC & AT & estag. & 22 \\
            DIMACS & HR & estag. & 27 \\
            DIMACS & HR & tempo & 23 \\
            DIMACS & AT & estag. & 47 \\
            DIMACS & AT & tempo & 3 \\
            HB & HR & iter. & 4 \\
            HB & HR & estag. & 25 \\
            HB & HR & tempo & 21 \\
            HB & AT & iter. & 7 \\
            HB & AT & estag. & 43 \\
            \bottomrule
        \end{tabular}
    }{
    \Fonte{Autor (2026).}
}
\end{table}
